\lysh{14122}{4}

\begin{alist}
\item 
\begin{rlist}
\item Αφού η περίμετρος του ορθογωνίου είναι ίση με $10\text{ cm}$, θα ισχύει:
\[2y + 2 \cdot (2x) = 10 \Leftrightarrow 2y + 4x = 10 \Leftrightarrow y + 2x = 5 \Leftrightarrow y = 5 - 2x.\]

Οι διαστάσεις του ορθογωνίου είναι θετικοί αριθμοί, οπότε:
\[\begin{cases} 2x > 0 \\ y > 0 \end{cases} \Leftrightarrow \begin{cases} x > 0 \\ 5 - 2x > 0 \end{cases} \Leftrightarrow \begin{cases} x > 0 \\ x < \dfrac{5}{2} \end{cases}\]
Τελικά: $x \in \left(0, \dfrac{5}{2}\right)$.
\item Το εμβαδόν (σε $\text{cm}^2$) του ορθογωνίου δίνεται από τη σχέση:
\[E_{\text{ορθ}} = y \cdot 2x = (5 - 2x) \cdot 2x = 10x - 4x^2, x \in \left(0, \dfrac{5}{2}\right).\]
\end{rlist}

\item Το μέρος του εμβαδού (σε $\text{cm}^2$) του ορθογωνίου που βρίσκεται έξω από τον κύκλο είναι:
\[E = E_{\text{ορθ}} - E_{\text{κύκλου}} = 10x - 4x^2 - \pi \cdot x^2 = 10x - (4 + \pi)x^2, x \in \left(0, \dfrac{5}{2}\right).\]

\item 
\begin{rlist}
\item Το εμβαδό $E$ του ορθογωνίου που βρίσκεται έξω από τον κύκλο πρέπει να είναι ίσο
με $(6 - \pi)\text{ cm}^2$, οπότε έχουμε:
\[E = 6 - \pi \Leftrightarrow 10x - (4 + \pi)x^2 = 6 - \pi \Leftrightarrow (4 + \pi)x^2 - 10x + (6 - \pi) = 0, \quad (1).\]

Η διακρίνουσα του τριωνύμου είναι:
\[\Delta = 100 - 4 \cdot (4 + \pi) \cdot (6 - \pi) = 4 - 8\pi + 4\pi^2 = 4(1 - 2\pi + \pi^2) = 4(1 - \pi)^2 > 0,\]
οπότε η εξίσωση $(1)$ έχει δυο λύσεις άνισες τις:
$x = \dfrac{10 + 2(1 - \pi)}{2(4 + \pi)} = \dfrac{12 - 2\pi}{2(4 + \pi)} = \dfrac{2(6 - \pi)}{2(4 + \pi)} = \dfrac{6 - \pi}{4 + \pi}$, που απορρίπτεται γιατί ο $x$ είναι ρητός.
και
\[x = \dfrac{10 - 2(1 - \pi)}{2(4 + \pi)} = \dfrac{8 + 2\pi}{2(4 + \pi)} = \dfrac{2(4 + \pi)}{2(4 + \pi)} = 1.\]
Τελικά η ακτίνα του κύκλου είναι $x = 1\text{ cm}$.

\item Οι διαστάσεις του ορθογωνίου είναι $2x = 2\text{ cm}$ και $y = 5 - 2 \cdot 1 = 3\text{ cm}$.
\end{rlist}
\end{alist}
